\section{Introduction}

The \SF problem is one of the most well-studied problems in network
design~\cite{ljubic2021solving, hwang1992steiner, du2013advances,
gupta2011approximation}. In this problem the input consists of a graph
$G=(V,E)$ with positive edge weights, %
a set of \emph{terminals} $R\subseteq V$, and a set of
\emph{demands}~$D\subseteq\binom{R}{2}$. The objective is to select a subgraph
$F\subseteq G$, minimizing the total cost of selected edges, while ensuring that for every
demand pair $\{s,t\}\in D$, $s$ and $t$ are in the same connected component of
$F$.  Since edge weights are positive, it is easy to see that the optimal
solution is always a forest. The \SF problem finds many applications (see
surveys~\cite{ljubic2021solving, cheng2013steiner, voss2006steiner,
tang2020survey}), for example in telecommunication networks
(cf.~\cite{voss2006steiner}). 

Our goal in this paper is to reassess the complexity of this fundamental problem
from the point of view of parameterized complexity and approximation
algorithms.\footnote{We assume the reader is familiar with the
basics of parameterized complexity and approximation algorithms, such as the classes FPT and APX
and the definition of treewidth, as given in standard textbooks 
\cite{cygan2015parameterized,williamson2011design,feldmann2020survey}. 
We give full definitions of all parameters in \cref{sec:prelims}.}
In order to recall the context, it is helpful to compare \SF to the
even more well-studied \ST problem, which is the special case of \SF where all
terminals are required to be connected, i.e., $D=\binom{R}{2}$, and an optimal
solution is a tree. \ST was already included in Karp's seminal
list~\cite{karp1975computational} of NP-hard problems from the 1970s. From the
approximation point of view, \ST (and therefore \SF) is known to be
APX-hard~\cite{chlebik2008steiner}, but both problems admit constant factor
approximations in polynomial time for general input graphs, where the best
approximation factors known are~$\ln(4)+\eps<1.39$~\cite{byrka2013steiner}
and~$2$~\cite{agrawal1991trees, ravi1994primal}, respectively. Despite this
similarity, when considering graph width parameters the problems exhibit 
wildly divergent behaviors from the parameterized complexity point of
view: whereas \ST is FPT parameterized by standard structural parameters such
as treewidth and can in fact even be solved in single exponential
$2^{O(k)}n^{O(1)}$ time~\cite{bodlaender2015deterministic} when $k$ is the
treewidth, \SF is NP-hard on graphs of treewidth~$3$, as shown independently by
\citet{gassner2010steiner} and \citet{bateni2011approximation}.

\SF is therefore a problem that presents a dramatic jump in complexity in this
context, compared to \ST, as the hardness result on graphs of treewidth~$3$
rules out even an XP algorithm for parameter treewidth. One of the main
positive contributions of \citet{bateni2011approximation} was an algorithm
attempting to bridge this gap using approximation. In particular, they showed
that \SF admits an approximation scheme for graphs of treewidth~$k$, which
computes a $(1+\eps)$\hy{}approximation in $n^{O(k^2/\eps)}$ time for any $\eps>0$. Hence, if we
allow slightly sub-optimal solutions, we can at least place the problem in XP
parameterized by treewidth. In their paper, \citet{bateni2011approximation}
remark that because the exponent of the polynomial of this runtime depends on
$k$ and $\eps$, "it remains an interesting question for future research whether
this dependence can be removed", that is, whether a $(1+\eps)$-approximation
can be obtained in FPT time.

The main result of our paper is a positive resolution of the question of
\citet{bateni2011approximation}: we show that \SF\ admits an \emph{efficient
parameterized approximation scheme (EPAS)} for treewidth, that is, a
$(1+\eps)$-approximation algorithm with a runtime of the form
$f(k,\eps)n^{O(1)}$. In other words, we show that their algorithm can be
improved in a way that makes the running time FPT not only in the treewidth,
but also in $1/\eps$. More precisely, we show the following:



\begin{theorem}\label{thm:tw-EPAS}
The \SF problem admits an EPAS parameterized by the treewidth $k$ with a runtime of $2^{O(\frac{k^2}{\eps}\log\frac{k}{\eps})}\cdot n^{O(1)}$.
\end{theorem}


%

Moving on from treewidth, we ask what the most general
parameter is for which we may hope to obtain an FPT \emph{exact} algorithm for
\SF. We observe that the NP-hardness result of \citet{gassner2010steiner,
bateni2011approximation} for \SF on graphs of treewidth~$3$ actually has some
further implications for some even more restricted parameters: the graphs
constructed in their reductions also have constant \emph{treedepth} and
\emph{feedback vertex set} size, implying that the problem remains hard for
both of these parameters (which are incomparable in general). More
precisely, known reductions imply the following:

\begin{theorem}[\citet{gassner2010steiner,bateni2011approximation}]
The \SF problem is NP-hard on graphs of treewidth~$3$, treedepth~$4$, and feedback vertex set of size~$2$.
\end{theorem}

This leads us to consider even more restricted parameters, such as the size of
a \emph{vertex cover} and \emph{feedback edge set}, which are not bounded in
this reduction. Indeed, not only do we prove that \SF is FPT for both of these
parameters, but we are also able to determine the correct parameter dependence,
under the Exponential Time Hypothesis (ETH). For feedback edge set the optimal dependence is single
exponential:  

\begin{theorem}\label{thm:fes}
The \SF problem is FPT parameterized by the size $k$ of a feedback edge set and can be solved in $2^{O(k)}n^{O(1)}$ time.
Furthermore, no $2^{o(k)}n^{O(1)}$ time algorithm exists, under ETH.
\end{theorem}

For the parameterization by the vertex cover size, we obtain a slower runtime
for our FPT algorithm. Interestingly, we are also able to prove that this is
best possible, under~ETH. Our lower bound for \SF is in contrast to the \ST
problem, for which a faster $2^{O(k)}n^{O(1)}$ time algorithm exists, even if
$k$ is the treewidth~\cite{bodlaender2015deterministic}.

\begin{theorem}\label{thm:vc}
The \SF problem is FPT parameterized by the size $k$ of a vertex cover and can be solved in $2^{O(k\log k)}n^{O(1)}$ time. Furthermore, no $2^{o(k\log k)}n^{O(1)}$ time algorithm exists, under ETH.
\end{theorem}

We remark that \citet{bodlaender2023complexity} recently independently showed 
that \SF admits a~$2^{O(k\log k)}n^{O(1)}$ time algorithm for the size $k$ of a 
vertex cover (improving an algorithm for the unweighted version of the problem 
given in~\cite{GimaHKKO22}). While they develop their own dynamic program to 
solve this problem, we rely on an existing algorithm 
by~\cite{bateni2011approximation} (see \autoref{thm:tw-dp}). Accordingly our 
description of the algorithm is very short compared 
to~\cite{bodlaender2023complexity}. The more interesting part of 
\autoref{thm:vc} however is the proof of the lower bound.

%
%
%

\subsection{Overview of Techniques}

%
%
%
%
Let us briefly sketch the high level ideas of our results given by \cref{thm:tw-EPAS,thm:vc,thm:fes}.

\smallskip
\textbf{EPAS for treewidth.} Our algorithm extends
the work of \cite{bateni2011approximation}, so let us briefly recall some key
ideas. 
Given a rooted tree decomposition, a terminal $t$ is called \emph{active} for a bag $B$
if there is a demand $\{s,t\}\in D$ such that $t$ lies in the sub-tree rooted at $B$
while $s$ does not (see \cref{sec:prelims} for formal definitions). 
It is a standard property of tree decompositions that every bag is a separator.
Hence the
component of any feasible solution that contains an active terminal must intersect~$B$. 
The hardness of the problem now inherently stems from the fact
that we have to decide for all active terminals of a bag, how
the corresponding component intersects the bag, and therefore how the active
terminals (whose number is unbounded by $k$) are partitioned into connected
components.  Suppose, however, that someone supplied us with this information,
that is, suppose that for each bag $B$ we are given a set of partitions $\Pi_B$
of its active terminals and we are promised that the optimal solution
\emph{conforms} to all $\Pi_B$. By this we mean that if we look at how the
optimal solution partitions the active terminals of $B$ into connected
components and call this partition $\pi$, then $\pi\in \Pi_B$, that is, the
optimal partition is always one of the supplied options. In this case, using
this extra information, the problem does become tractable, as shown in~\cite{bateni2011approximation}:

\begin{theorem}[\citet{bateni2011approximation}]\label{thm:tw-dp} For an input
graph $G$ on $n$ vertices, let a rooted nice tree decomposition of width~$k$ be
given, such that all terminals lie in bags of leaf nodes of the decomposition.
Also, let a set~$\Pi_B$ of partitions of the active terminals of each bag $B$ of
the decomposition be given. If $p=\sum_B |\Pi_B|$ is the total number of
partitions, then a minimum cost \SF solution conforming to all $\Pi_B$ can be
computed in $2^{O(k\log k)}\cdot(pn)^{O(1)}$ time.  
\end{theorem}

The above theorem does not seem immediately helpful, since one would still need to
find a small collection of partitions $\Pi_B$ in order to obtain an efficient algorithm. 
Note however, that the partition sets may conform to an approximate solution as well, which
would let the algorithm compute a solution that it at least as good.
The strategy of \cite{bateni2011approximation} therefore is to
construct a collection of partitions that has size polynomial in~$n$ (when $k,\eps$ are
fixed constants) by stipulating that when two active terminals are ``close'' to
each other, they should belong in the same set of the partition of some near-optimal solution. In order to
bound the resulting approximation ratio, they need to provide a
charging scheme: starting from an optimal solution, they merge components which are
``close'', to obtain a solution that conforms to the~$\Pi_B$ used by the algorithm. 
They then show that the resulting solution is still near-optimal by charging the extra
cost incurred by a merging operation to one of the two merged components.

A blocking point in the above is that we need to make sure we do not
``overcharge'' any component. This is accomplished in~\cite{bateni2011approximation}
via a partial ordering of the components: we
order the components according to the highest bag of the rooted tree decomposition they
intersect, and whenever two components are merged we charge this to the
\emph{lower} component. As shown in~\cite{bateni2011approximation},
this ensures that no component is charged for more than
$k$ merges. Unfortunately, this also implies that the merging procedure is not
symmetric, which severely diminishes the contexts in which we can apply it.

Let us now describe how our approach improves upon this algorithm. A key
ingredient  will be a more sophisticated charging scheme, which will allow us
to obtain a better (smaller) collection of partitions~$\Pi_B$, without
sacrificing solution quality.  Counter-intuitively, we will achieve this by
introducing a \emph{second} parameter: the height $h$ of the tree
decomposition.  Informally, we will now construct a near-optimal solution by
merging two components whenever the connection cost is low compared to the cost of
(a~part~of) \emph{either} component (as opposed to the lower component). As 
in~\cite{bateni2011approximation}, this runs the risk of charging many merging
operations to a higher component, but by performing an accounting by tree
decomposition level and using the fact that the decomposition only has $h$
levels, we are able to show that our solution is still near-optimal even though
we merge components much more aggressively than~\cite{bateni2011approximation}.
In this way, for each bag we construct one partition of its active terminals into
a number of sets that is polynomial in $k+h+\frac{1}{\eps}+\log n$, in a way
that guarantees that this partition is a \emph{refinement} of a near-optimal
solution. That is, whenever we decide to place two terminals together in our
partition, the near-optimal solution does the same. However, this solution does
not necessarily conform to the resulting partitions, as two terminals of the same component
might end up in different sets of the partition for a bag.

At this point an astute reader may be wondering that since we consider both the
width $k$ and the height $h$ of the decomposition as parameters, we are
effectively parameterizing by treedepth, rather than treewidth. This is
correct, but we then go on to invoke a result of \citet{BodlaenderH98} which
states that any tree decomposition can be rebalanced to have height $O(\log n)$
without severely increasing its width. Hence, the family of partitions we now
have has size polynomial in $k+\frac{1}{\eps}+\log n$. 
However, we are not done yet, since at this point we can
only guarantee that our partitions are refinements of a near-optimal
solution. To complete the algorithm, we work from this family of partitions to
obtain a collection of partitions \emph{conforming} to our near-optimal
solutions using $\delta$-nets (this is similar to the approach of
\cite{bateni2011approximation}). This leads to a running time of the form
$(\log n)^{O(\frac{k^2}{\eps})}n^{O(1)}$, which by standard arguments of parameterized
complexity is in fact FPT and can be upper-bounded by a function of the form
$2^{O(\frac{k^2}{\eps}\log\frac{k}{\eps})}\cdot n^{O(1)}$. 
To summarize, our high-level strategy is to
show that the approach of \cite{bateni2011approximation} can be significantly
improved when the input decomposition has small width and height, but then we
observe that our new scheme is efficient enough in the height that even if we
replace $h$ by a bound that can be obtained for \emph{any} graph, we still have
an algorithm with an FPT running time, that is, significantly faster than that
of \cite{bateni2011approximation}. 

\smallskip
\textbf{Vertex Cover.} For the parameterization by the vertex cover size, as
mentioned we obtain an FPT exact algorithm with dependence $2^{O(k\log k)}$. A
similar algorithm was recently independently obtained by
\citet{bodlaender2023complexity} via dynamic programming. However, our
algorithm is significantly simpler, because our strategy is to show how to
construct a tree decomposition and a collection of partitions $\Pi_B$ such that
we only need \emph{one} partition of the active terminals for each bag. As a
consequence, $p=O(n)$ and \autoref{thm:tw-dp} implies the algorithm of
\autoref{thm:vc}, without the need to formulate a new dynamic program.

Our main result for this parameter is that under ETH the runtime dependence
is asymptotically optimal. Note that this also implies that the runtime of the dynamic program
given by \autoref{thm:tw-dp} cannot be improved with regards to the dependence
on the treewidth. To show this, we present a reduction from \textsc{3-SAT}, where the
goal is to compress an $n$-variable formula into a \SF
instance such that the graph has vertex cover size $O(n/\log n)$. The intuition on why
it is possible to achieve such a compression is the following: suppose we have
an instance with vertex cover of size $k$ and a demand between two vertices of
the independent set. Then the simplest way to satisfy such a demand is to
connect both vertices to a common neighbor in the vertex cover. This encodes a
choice among $k$ vertices, and hence it is sufficient to encode the assignment for
$\log(k)$ binary variables. The strategy of our reduction is to set up some
choice gadgets which allow us to encode the assignments to the original formula
taking advantage of the fact that each choice can represent a logarithmic
number of variables. Hence we can obtain a construction of slightly sub-linear
($O(n/\log n)$) size.  We then of course need to add some verification gadgets,
representing the clauses, to check that the formula is indeed satisfied. But
even though the number of such gadgets is linear in $n$, we make sure that they
form an independent set, and hence the total vertex cover size remains
sufficiently small to obtain our lower bound. We note that this compression
strategy is similar to techniques recently used to obtain slightly
super-exponential lower bounds for vertex cover for other problems
\cite{LampisMV23,LampisV23}, but the constructions we use are new and tailored to
\SF.

\smallskip
\textbf{Feedback Edge Set.} For the parameterization by the size $k$ of a
feedback edge set, instead of relying on the dynamic program given by
\autoref{thm:tw-dp} we go an entirely different route in order to obtain the
faster $2^{O(k)}n^{O(1)}$ time FPT algorithm of \autoref{thm:fes}. First off,
it is not hard to reduce the \SF problem to an instance in which all vertices
have degree at least~$2$.  We then consider paths with internal vertices of
degree $2$, with endpoints that are vertices incident to the feedback edge set
or vertices of degree at least~$3$.  We call these paths \emph{topo-edges} and
argue that there are only $O(k)$ of these. We then guess for which topo-edges
the two endpoints lie in different components of the optimal \SF solution,
which can be done in $2^{O(k)}$ time.  If a topo-edge has both its endpoints in
the same component of the optimum, we show that it can be easily handled. For
the remaining topo-edges, we can decide which edges along the path do not
belong to the optimal solution by a reduction to the polynomial-time solvable
\pname{Min Cut} problem.



%
%
%
%
%
%
%
%
%
%
%
%
%
%
%
%
%
%
%
%
%
%
%
%
%
%
%
%
%
%


%
%





\subsection{Related work}
\label{sec:related}

\citet{bateni2011approximation} show that one of the consequences of their XP
approximation scheme is a PTAS for \SF on planar graphs, by using the common
technique pioneered by \citet{baker1994approximation} of reducing this problem
to graphs for which the treewidth is bounded as a function of~$\eps$.  Because
their algorithm is not FPT, their PTAS has a running time of the form
$n^{f(\eps)}$. By using our algorithm from \autoref{thm:tw-EPAS} we can improve
this runtime to~$f(\eps)n^{O(1)}$, i.e., we obtain on EPTAS for planar graphs.
However, \citet{eisenstat2012efficient} already showed that a
$(1+\eps)$-approximation algorithm with a runtime of $O(f(\eps)\cdot n\log^3
n)$ exists for \SF on planar graphs. While they build on the work of
\citet{bateni2011approximation}, and in particular also reduce to graphs of
treewidth bounded as a function of~$\eps$, interestingly they do not obtain an
EPAS parameterized by treewidth.  Instead they use a different route and show
that given a graph~$H$ of treewidth~$k$, in $O(f(k,\eps)\cdot n\log^2 n)$ time
it is possible to compute a \SF solution in $H$ whose cost is at most
$\cost(F^\star)+\eps\cost(H)$, i.e., there is an additive error that depends on
the cost of~$H$ compared to the optimum solution~$F^\star$. If the input
graph~$G$ is planar, then a result by \citet{borradaile2009n} implies that from
$G$ a so-called \emph{banyan}~\cite{rao1998approximating, bartal2021near} can
be computed, which is a subgraph of $G$ with cost bounded by
$O(g(\eps)\cost(F^\star))$, and which contains a near-optimal approximation of
every Steiner forest (cf.~\cite[Lemma~2.1]{eisenstat2012efficient}). By
applying the framework of \citet{bateni2011approximation} on the banyan instead
of the input graph, it is then possible to obtain a graph $H$ of treewidth
bounded by a function of~$\eps$, for which the algorithm of
\citet{eisenstat2012efficient} computes a $(1+O(\eps))$-approximation for the
input. 

If it would be possible to compute a banyan for bounded treewidth graphs, then the algorithm of \citet{eisenstat2012efficient} would also imply an EPAS for treewidth. However, to the best of our knowledge, and as explicitly stated by \citet{bartal2021near}, banyans are only known for planar graphs~\cite{eisenstat2012efficient, borradaile2009n}, Euclidean metrics~\cite{rao1998approximating}, and doubling metrics~\cite{bartal2021near} (in fact, the latter are so-called \emph{forest banyans}, which have weaker properties). Thus it is unclear how to obtain an EPAS for \SF parameterized by the treewidth via the algorithm of \citet{eisenstat2012efficient}. We leave open whether a banyan exists for bounded treewidth graphs, which could give an alternative algorithm to the one given in \autoref{thm:tw-EPAS}. However, a further remark is that the cost of the banyan for planar graphs obtained by \citet{borradaile2009n} has exponential dependence on~$1/\eps$, which implies a double exponential runtime dependence on~$1/\eps$ for the EPTAS for planar graphs. If a banyan can be obtained for bounded treewidth graphs by generalizing the techniques of \citet{borradaile2009n} to minor-free graphs, then the resulting EPAS parameterized by treewidth would also have double exponential runtime in~$1/\eps$. In this case however, our EPAS given by \autoref{thm:tw-EPAS} would be exponentially faster.

A different parameter that is often studied in the context of Steiner problems 
is the number $p=|R|$ of terminals. The classic result of 
\citet{dreyfus1971steiner} presents an FPT algorithm for \ST with a runtime of 
$3^p n^{O(1)}$. For unweighted graphs, this was 
improved~\cite{bjorklund2007fourier, nederlof2009fast} to $2^p n^{O(1)}$, while 
the fastest known algorithm for weighted graphs can compute the optimum in 
$(2+\eps)^p n^{O(\sqrt{\frac{1}{\eps}}\log\frac{1}{\eps})}$ 
time~\cite{fuchs2007dynamic} for any~$\eps>0$. The algorithm of 
\citet{dreyfus1971steiner} can be generalized to solve \SF in $2^{O(p)} 
n^{O(1)}$ time (cf.~\cite{chitnis2021parameterized}). A somewhat dual parameter 
to the number of terminals is the number~$q$ of non-terminals (so-called 
\emph{Steiner vertices}) in the optimum solution. For this parameter, a folklore 
result states that \ST (and thus also \SF) is 
W[2]-hard~(cf.~\cite{cygan2015parameterized, dvorak2021parameterized}). However, 
an EPAS with a runtime of $2^{O(q^2/\eps^4)}n^{O(1)}$ was shown to exist for 
\ST~\cite{dvorak2021parameterized}. For \SF it is not hard to see that such an 
EPAS parameterized by $q$ cannot exist unless P=NP 
(cf.~\cite{dvorak2021parameterized}), but if $c$ denotes the number of 
components of the optimum solution, there is an EPAS with a runtime of 
$(2c)^{O((q+c)^2/\eps^4)}n^{O(1)}$~\cite{dvorak2021parameterized}. Similar 
results have been found for related Steiner problems in directed 
graphs~\cite{chitnis2021parameterized}. For further results in the area of 
parameterized approximations, we refer to the survey~\cite{feldmann2020survey}.
